\documentclass[11pt]{article}

\usepackage{amsmath,amssymb,amsthm,mathtools}
\usepackage[margin=1in]{geometry}

\newtheorem{theorem}{Theorem}
\newtheorem{lemma}[theorem]{Lemma}
\newtheorem{corollary}[theorem]{Corollary}
\newtheorem{remark}[theorem]{Remark}
\newtheorem{definition}[theorem]{Definition}

\begin{document}

\section{Closure under vertex and edge self-amalgamation}

For a finite graph $H$ and $0\le p<1$, define
\[
    \Phi_H(p)
    :=
    (1-p)^{v(H)}
    \chi_H\!\left(\frac{1}{1-p}\right),
\]
with the continuous extension at $p=1$.

Since $\chi_H$ is monic of degree $v(H)$, this extension has value
$\Phi_H(1)=1$.  Throughout, a graphon means a symmetric, jointly
measurable function $W:\Omega^2\to[0,1]$ on an arbitrary probability
space $(\Omega,\mathcal A,\mu)$.  We write
\[
    p=t(K_2,W).
\]

\subsection{Definitions}

\begin{definition}[Vertex self-amalgam]
Let $F$ be a finite graph with a distinguished vertex $a\in V(F)$,
and let $N\ge1$.

Take $N$ disjoint labelled copies
\[
    F_1,\dots,F_N
\]
of the rooted graph $(F,a)$ and identify all $N$ copies of the
distinguished vertex $a$, making no other identifications.

The resulting graph is denoted by
\[
    B_N^{(1)}(F,a)
\]
and is called the $N$-fold vertex self-amalgam of $(F,a)$.
\end{definition}

\begin{definition}[Edge self-amalgam]
Let $F$ be a finite graph with a distinguished ordered edge
\[
    (a,b),
    \qquad ab\in E(F),
\]
and let $N\ge1$.

Take $N$ disjoint labelled copies of the rooted graph $(F,a,b)$.
Identify all copies of $a$ with one another and all copies of $b$
with one another, respecting the order of the distinguished edge,
and make no other identifications.

The resulting graph is denoted by
\[
    B_N^{(2)}(F,a,b)
\]
and is called the $N$-fold edge self-amalgam of $(F,a,b)$.
\end{definition}

The ordering of the distinguished edge is only used to specify
unambiguously how the copies are glued. It is important when the two
endpoints of the distinguished edge are not interchangeable by an
automorphism of the rooted graph.

\subsection{The graphon gluing inequalities}

We first prove the analytic part directly in graphon space.

\begin{lemma}[Vertex self-amalgamation inequality]
Let $F$ be rooted at a vertex $a$, and let $N\ge1$. Then for every
graphon $W$ on every probability space,
\[
    t(B_N^{(1)}(F,a),W)
    \ge
    t(F,W)^N.
\]
\end{lemma}

\begin{proof}
For $x\in\Omega$, define the rooted homomorphism density
\[
    f(x)
    :=
    \int_{\Omega^{V(F)\setminus\{a\}}}
       \prod_{uv\in E(F)}
       W(z_u,z_v)
       \prod_{v\in V(F)\setminus\{a\}}d\mu(z_v),
\]
where
\[
    z_a=x.
\]
By Fubini's theorem,
\[
    t(F,W)=\int_\Omega f(x)\,d\mu(x).
\]

In the graph $B_N^{(1)}(F,a)$, the $N$ copies of
$F\setminus\{a\}$ are pairwise vertex-disjoint once the image of the
common root is fixed. Hence their integrations factor, giving
\[
    t(B_N^{(1)}(F,a),W)
    =
    \int_\Omega f(x)^N\,d\mu(x).
\]
Since $f\ge0$ and $u\mapsto u^N$ is convex on $[0,\infty)$,
Jensen's inequality gives
\[
\begin{aligned}
    t(B_N^{(1)}(F,a),W)
    &=
    \int_\Omega f(x)^N\,d\mu(x)
    \\
    &\ge
    \left(\int_\Omega f(x)\,d\mu(x)\right)^N
    \\
    &=
    t(F,W)^N.
\end{aligned}
\]
All integrands are measurable and bounded between $0$ and $1$.
Thus Tonelli--Fubini applies, including on probability spaces that are
not atomless or standard Borel.
\end{proof}

\begin{lemma}[Edge self-amalgamation inequality]
Let $F$ have a distinguished ordered edge $(a,b)$, let $N\ge1$, and
let $W$ be a graphon on an arbitrary probability space with
\[
    p=t(K_2,W)>0.
\]
Then
\[
    t(B_N^{(2)}(F,a,b),W)
    \ge
    \frac{t(F,W)^N}{p^{N-1}}.
\]
\end{lemma}

\begin{proof}
For $(x,y)\in\Omega^2$, define
\[
    g(x,y)
    :=
    \int_{\Omega^{V(F)\setminus\{a,b\}}}
       \prod_{uv\in E(F)\setminus\{ab\}}
       W(z_u,z_v)
       \prod_{v\in V(F)\setminus\{a,b\}}d\mu(z_v),
\]
where
\[
    z_a=x,
    \qquad
    z_b=y.
\]
The distinguished edge $ab$ has deliberately been omitted from the
product defining $g$. Therefore
\begin{equation}\label{eq:tF-edge-rooted}
    t(F,W)
    =
    \int_{\Omega^2}
       W(x,y)g(x,y)\,d\mu(x)d\mu(y).
\end{equation}

In the $N$-fold edge self-amalgam the distinguished edge is a single
edge of the resulting simple graph, rather than $N$ parallel edges.
Once the images $x,y$ of the two common root vertices have been fixed,
all remaining vertices in the $N$ copies are independent. Consequently,
\begin{equation}\label{eq:edge-amalgam-density}
    t(B_N^{(2)}(F,a,b),W)
    =
    \int_{\Omega^2}
       W(x,y)g(x,y)^N\,d\mu(x)d\mu(y).
\end{equation}

Since
\[
    p=\int_{\Omega^2}W(x,y)\,d\mu(x)d\mu(y)>0,
\]
the formula
\[
    d\rho(x,y)
    :=
    \frac{W(x,y)}{p}\,d(\mu\otimes\mu)(x,y)
\]
defines a probability measure on $\Omega^2$.
Equations \eqref{eq:tF-edge-rooted} and
\eqref{eq:edge-amalgam-density} become
\[
    t(F,W)
    =
    p\int g\,d\rho
\]
and
\[
    t(B_N^{(2)}(F,a,b),W)
    =
    p\int g^N\,d\rho.
\]
Applying Jensen's inequality to $g\ge0$ gives
\[
\begin{aligned}
    t(B_N^{(2)}(F,a,b),W)
    &=
    p\int g^N\,d\rho
    \\
    &\ge
    p\left(\int g\,d\rho\right)^N
    \\
    &=
    p\left(\frac{t(F,W)}{p}\right)^N
    \\
    &=
    \frac{t(F,W)^N}{p^{N-1}}.
\end{aligned}
\]
As above, the finite products are measurable and bounded, so all
factorizations and changes in integration order are justified by
Tonelli--Fubini.
\end{proof}

\subsection{Chromatic-polynomial identities}

The preceding analytic inequalities exactly match the
chromatic-polynomial target because the densities of $K_1$ and $K_2$
are completely determined by the edge density.

\begin{lemma}[Chromatic polynomial of a vertex self-amalgam]
Let
\[
    H=B_N^{(1)}(F,a).
\]
Then
\[
    \chi_H(x)
    =
    \frac{\chi_F(x)^N}{x^{N-1}}.
\]
Consequently,
\[
    \Phi_H(p)=\Phi_F(p)^N.
\]
\end{lemma}

\begin{proof}
Fix an integer $k\ge1$. By symmetry under permutations of the $k$
colours, the number of proper $k$-colourings of $F$ in which the
distinguished vertex $a$ is assigned any prescribed colour is
\[
    \frac{\chi_F(k)}{k}.
\]
To colour $H$, first choose the colour of the common root, and then
colour the remaining vertices in each of the $N$ copies independently.
Thus
\[
\begin{aligned}
    \chi_H(k)
    &=
    k
    \left(\frac{\chi_F(k)}{k}\right)^N
    \\
    &=
    \frac{\chi_F(k)^N}{k^{N-1}}.
\end{aligned}
\]
Because $F$ is nonempty, $\chi_F(0)=0$, so the factor theorem shows
that $x$ divides $\chi_F(x)$.  Hence the quotient below is a
polynomial.  The counting identity holds for every positive
integer $k$, so equality at infinitely many values gives the polynomial
identity
\[
    \chi_H(x)
    =
    \frac{\chi_F(x)^N}{x^{N-1}}.
\]

Let
\[
    n=v(F),
    \qquad
    q=1-p.
\]
Since the $N$ copies share exactly one vertex,
\[
    v(H)
    =
    Nn-(N-1).
\]
Therefore, for $p<1$,
\[
\begin{aligned}
    \Phi_H(p)
    &=
    q^{Nn-(N-1)}
    \frac{\chi_F(q^{-1})^N}
         {(q^{-1})^{N-1}}
    \\
    &=
    q^{Nn}
    \chi_F(q^{-1})^N
    \\
    &=
    \bigl(q^n\chi_F(q^{-1})\bigr)^N
    \\
    &=
    \Phi_F(p)^N.
\end{aligned}
\]
The identity at $p=1$ follows by continuity; both sides equal $1$.
\end{proof}

\begin{lemma}[Chromatic polynomial of an edge self-amalgam]
Let
\[
    H=B_N^{(2)}(F,a,b).
\]
Then
\[
    \chi_H(x)
    =
    \frac{\chi_F(x)^N}{(x)_2^{N-1}},
\]
where
\[
    (x)_2=x(x-1).
\]
Consequently, for $p>0$ (including $p=1$ by continuous extension),
\[
    \Phi_H(p)
    =
    \frac{\Phi_F(p)^N}{p^{N-1}}.
\]
\end{lemma}

\begin{proof}
Fix an integer $k\ge2$. Since $a$ and $b$ are adjacent, every proper
colouring assigns them two distinct colours.

The permutation group of the $k$ colours acts transitively on the
ordered pairs of distinct colours. Hence, for any prescribed ordered
pair $(c_1,c_2)$ with $c_1\ne c_2$, the number of proper colourings
of $F$ satisfying
\[
    a\mapsto c_1,
    \qquad
    b\mapsto c_2
\]
is
\[
    \frac{\chi_F(k)}{(k)_2},
    \qquad
    (k)_2=k(k-1).
\]
There are $(k)_2$ choices for the ordered pair of colours on the
common root edge, after which the $N$ copies may be coloured
independently. Thus
\[
\begin{aligned}
    \chi_H(k)
    &=
    (k)_2
    \left(
      \frac{\chi_F(k)}{(k)_2}
    \right)^N
    \\
    &=
    \frac{\chi_F(k)^N}{(k)_2^{N-1}}.
\end{aligned}
\]
Because $ab$ is an edge, $\chi_F(0)=\chi_F(1)=0$.  The factor theorem
therefore shows that $(x)_2=x(x-1)$ divides $\chi_F(x)$.  Thus the
quotient below is a polynomial.  Since the
counting identity holds for every integer $k\ge2$, equality at
infinitely many values gives
\[
    \chi_H(x)
    =
    \frac{\chi_F(x)^N}{(x)_2^{N-1}}.
\]

Let
\[
    n=v(F),
    \qquad
    q=1-p.
\]
The $N$ copies share exactly two vertices, so
\[
    v(H)
    =
    Nn-2(N-1).
\]
Moreover,
\[
\begin{aligned}
    (q^{-1})_2
    &=
    q^{-1}(q^{-1}-1)
    \\
    &=
    \frac1q\frac{1-q}{q}
    \\
    &=
    \frac{p}{q^2}.
\end{aligned}
\]
For $0<p<1$, therefore
\[
\begin{aligned}
    \Phi_H(p)
    &=
    q^{Nn-2(N-1)}
    \frac{\chi_F(q^{-1})^N}
         {(q^{-1})_2^{N-1}}
    \\
    &=
    q^{Nn-2(N-1)}
    \chi_F(q^{-1})^N
    \left(\frac{q^2}{p}\right)^{N-1}
    \\
    &=
    \frac{
      q^{Nn}\chi_F(q^{-1})^N
    }{
      p^{N-1}
    }
    \\
    &=
    \frac{\Phi_F(p)^N}{p^{N-1}}.
\end{aligned}
\]
At $p=1$, the same identity follows by continuity, and both sides are
$1$.  No claim involving the displayed quotient is needed at $p=0$.
\end{proof}

\subsection{Closure theorem}

\begin{theorem}[Closure under vertex and edge self-amalgamation]
Let $F$ be a finite graph, let $W$ be a graphon, and put
\[
    p=t(K_2,W).
\]
Assume
\[
    t(F,W)\ge \Phi_F(p)
\]
and
\[
    \Phi_F(p)\ge0.
\]

Then the following hold.

\begin{enumerate}
    \item
    For every distinguished vertex $a\in V(F)$ and every $N\ge1$,
    \[
        t(B_N^{(1)}(F,a),W)
        \ge
        \Phi_{B_N^{(1)}(F,a)}(p).
    \]

    \item
    If $p>0$, then for every distinguished ordered edge
    $(a,b)$ of $F$ and every $N\ge1$,
    \[
        t(B_N^{(2)}(F,a,b),W)
        \ge
        \Phi_{B_N^{(2)}(F,a,b)}(p).
    \]
\end{enumerate}
\end{theorem}

\begin{proof}
For the vertex self-amalgam,
\[
\begin{aligned}
    t(B_N^{(1)}(F,a),W)
    &\ge
    t(F,W)^N
    \\
    &\ge
    \Phi_F(p)^N
    \\
    &=
    \Phi_{B_N^{(1)}(F,a)}(p).
\end{aligned}
\]
The second inequality uses
\[
    t(F,W)\ge\Phi_F(p)\ge0.
\]

For the edge self-amalgam,
\[
\begin{aligned}
    t(B_N^{(2)}(F,a,b),W)
    &\ge
    \frac{t(F,W)^N}{p^{N-1}}
    \\
    &\ge
    \frac{\Phi_F(p)^N}{p^{N-1}}
    \\
    &=
    \Phi_{B_N^{(2)}(F,a,b)}(p).
\end{aligned}
\]
Again the passage to the $N$th power is justified by
$\Phi_F(p)\ge0$.
\end{proof}

\begin{remark}[Why the nonnegativity assumption is needed]
The condition
\[
    \Phi_F(p)\ge0
\]
cannot simply be omitted from this abstract argument. From
\[
    t(F,W)\ge\Phi_F(p)
\]
with $\Phi_F(p)<0$, one cannot conclude
\[
    t(F,W)^N\ge\Phi_F(p)^N
\]
when $N$ is even.

In the applications relevant to the Goodman-type problem this causes
no difficulty. For example, the target for an odd cycle is nonnegative
for $p\ge1/2$, and the pure-chordal target is nonnegative throughout
the range in which the clique-ratio proof applies.
\end{remark}

\begin{corollary}[Preservation of the required chromatic threshold]
Let
\[
    r:=\chi(F)\ge3.
\]
Suppose that for every graphon $W$ with
\[
    p=t(K_2,W)\ge1-\frac1{r-1}
\]
one has
\[
    t(F,W)\ge\Phi_F(p)\ge0.
\]
Then every vertex self-amalgam
\[
    B_N^{(1)}(F,a)
\]
and every edge self-amalgam
\[
    B_N^{(2)}(F,a,b)
\]
satisfies
\[
    t(H,W)\ge\Phi_H(p)
\]
throughout the same range
\[
    p\ge1-\frac1{r-1}.
\]
Moreover,
\[
    \chi(H)=r
\]
for both kinds of self-amalgam.
\end{corollary}

\begin{proof}
Because $r\ge3$, the displayed density range has $p\ge1/2>0$.
The density inequalities therefore follow immediately from the preceding
theorem, including its edge-amalgamation part.
It remains only to verify the chromatic number.

Every self-amalgam contains a copy of $F$, and therefore
\[
    \chi(H)\ge\chi(F)=r.
\]

For a vertex self-amalgam, take a proper $r$-colouring of each copy.
By permuting the $r$ colour labels separately in each copy, we may
arrange that the distinguished vertex has the same colour in every
copy. These colourings then combine into a proper $r$-colouring of the
amalgam. Thus
\[
    \chi(H)\le r.
\]

For an edge self-amalgam, the two endpoints of the distinguished edge
receive distinct colours. The symmetric group on the $r$ colours is
transitive on ordered pairs of distinct colours. Hence, after
permuting the colour labels independently in each copy, we may arrange
that the ordered pair of colours on the distinguished edge is the
same in every copy. The colourings then combine to give a proper
$r$-colouring of the amalgam.

Thus in either case
\[
    \chi(H)=r.
\]
Consequently the threshold
\[
    p\ge1-\frac1{\chi(H)-1}
\]
is exactly the original threshold
\[
    p\ge1-\frac1{\chi(F)-1}.
\]
\end{proof}

\subsection{The six-vertex Atlas application}

\begin{corollary}[Atlas 115]\label{cor:atlas115}
NetworkX Atlas graph $115$, with graph6 code \texttt{E`Xg}, is positive:
for every graphon $W$ of edge density $p\in[1/2,1]$,
\[
    t(F_{115},W)\geq p^3(2p-1)^2
    =\Phi_{F_{115}}(p).
\]
Its order, size, and chromatic number are respectively $6$, $7$, and $3$,
and
\[
    \chi_{F_{115}}(x)=x(x-1)^3(x-2)^2.
\]
\end{corollary}

\begin{proof}
Atlas 15 is the paw, with edge list
\[
  03,\ 12,\ 13,\ 23,
\]
where $123$ is its triangle and $0$ is the leaf at $3$.  Its
chromatic polynomial and target are
\[
  \chi_{F_{15}}(x)=x(x-1)^2(x-2),
  \qquad
  \Phi_{F_{15}}(p)=p^2(2p-1).
\]
The rooted triangle--tree theorem (the case $L_1$) proves
$t(F_{15},W)\geq\Phi_{F_{15}}(p)$ for every $p\in[1/2,1]$; the target
is nonnegative on this interval.

Take two copies of the paw and identify, in order, the endpoints of the
triangle edge $12$, the edge opposite the leaf-bearing vertex $3$.
The result is $B_2^{(2)}(F_{15},1,2)$.  One convenient edge labelling of
this amalgam is
\[
  01,\ 03,\ 05,\ 13,\ 15,\ 23,\ 45.
\]
It is isomorphic to Atlas 115, whose NetworkX edge list is
\[
  01,\ 14,\ 15,\ 23,\ 24,\ 25,\ 45.
\]
For example, the first list is carried to the second by
\[
  0\mapsto4,\quad 1\mapsto5,\quad 2\mapsto0,\quad
  3\mapsto1,\quad 4\mapsto3,\quad 5\mapsto2.
\]
The edge-amalgamation identities give
\[
  \chi_{F_{115}}(x)
  =\frac{\chi_{F_{15}}(x)^2}{x(x-1)}
  =x(x-1)^3(x-2)^2
\]
and, since $p\geq1/2>0$,
\[
  t(F_{115},W)
  \geq\frac{t(F_{15},W)^2}{p}
  \geq\frac{\Phi_{F_{15}}(p)^2}{p}
  =p^3(2p-1)^2.
\]
The amalgam contains the paw and is $3$-colourable by matching the
ordered colours on the common edge, so its chromatic number is $3$.
At $p=1/2$ the target is $0$; at $p=1$ both the target and the density
are $1$.  This proves the complete required interval.
\end{proof}

The other two edge orbits of the paw give Atlas 111 and Atlas 117:
gluing on a triangle edge incident with the leaf-bearing vertex gives
the former, while gluing on the pendant edge gives the latter.  Both
were already positive by the cone-over-forest theorem.  Exact
enumeration of the scoped Atlas graphs with at most six vertices shows
that Atlas 115 is the only formerly uncovered row added by vertex or
edge self-amalgamation of the accepted smaller positive bases; the
explicit isomorphism above supplies the proof-relevant identification.

\begin{remark}[Why only $K_1$ and $K_2$ give an automatic closure]
There is a useful reason that vertex and edge amalgamation are
particularly clean.

Suppose instead that copies of $F$ are amalgamated along a rooted
clique $K_s$. Provided $t(K_s,W)>0$, the same normalized-measure
argument gives
\[
    t(H,W)
    \ge
    \frac{t(F,W)^N}{t(K_s,W)^{N-1}}.
\]
On the chromatic-polynomial side, whenever $\Phi_{K_s}(p)>0$, one obtains
\[
    \Phi_H(p)
    =
    \frac{\Phi_F(p)^N}{\Phi_{K_s}(p)^{N-1}}.
\]
For $s=1$ and $s=2$,
\[
    t(K_1,W)=1=\Phi_{K_1}(p),
    \qquad
    t(K_2,W)=p=\Phi_{K_2}(p),
\]
so the analytic and chromatic denominators agree exactly.

For $s\ge3$, however, $t(K_s,W)$ is not determined by the edge
density $p$. Consequently the same argument does not by itself imply
the Goodman-type target. This is precisely where additional
clique-density or clique-ratio information is required.
\end{remark}

\subsection{Inputs, endpoint audit, and exact limitations}

The closure theorem itself uses only Tonelli--Fubini, Jensen's
inequality on a probability space, and the elementary colouring counts
above.  Corollary~\ref{cor:atlas115} additionally uses exactly the
accepted rooted triangle--tree bound for the paw $L_1$; it does not use
universal-vertex lifting, a finite-step reduction, or a regularity
assumption.

The division by $p^{N-1}$ in edge amalgamation is made only when
$p>0$.  This is automatic in the catalogue application and in the
required threshold corollary because $\chi(F)\ge3$.  The endpoint
$p=1$ is not obtained by substituting $q^{-1}$ with $q=0$: the target
identities there follow from their polynomial continuations, while the
graphon inequality itself follows directly from Jensen.  Vanishing
rooted densities cause no division, since the normalized measure
$d\rho=W\,d(\mu\otimes\mu)/p$ depends only on $p>0$.

Equality in the analytic vertex-amalgamation inequality occurs precisely
when the rooted density $f$ is constant almost everywhere (apart from
the tautological case $N=1$).  Equality in the edge version occurs when
$g$ is constant $\rho$-almost everywhere.  These observations are not
needed for positivity but identify the only possible slack in the
closure step.

Finally, the theorem applies only to copies that meet in exactly the
specified root vertex or root edge and have no other vertices or edges
between different copies.  It proves neither amalgamation along a
larger subgraph nor an operation in which different root types are mixed.
The clique discussion above explains why the same one-line comparison
does not automatically extend to $K_s$ for $s\ge3$.

%%%%%%%%%%%%%%%%%%%%%%%%%%%%%%%%%%%%%%%%%%%%%%%%%%%%%%%%%%%%%%%%%%%%%%%%%%%%%%%
\section{What the Lean formalization actually proves}
%%%%%%%%%%%%%%%%%%%%%%%%%%%%%%%%%%%%%%%%%%%%%%%%%%%%%%%%%%%%%%%%%%%%%%%%%%%%%%%

The machine-checked development is
\texttt{lean/Taeyoung/Methods/SelfAmalgam.lean}, ending at
\texttt{satisfiesLowerBound\_115}.

\paragraph{Scope: the single instance, not the closure theorem.}
Corollary~\ref{cor:atlas115} is proved in full, on all of $p\in[1/2,1]$.  The
general statements are \emph{not} formalized: neither the vertex
self-amalgamation inequality, nor the edge one at general $N$, nor the two
chromatic-polynomial identities, nor the closure theorem, nor the
threshold-preservation corollary.  What is formalized is the $N=2$ edge
amalgam of the paw, which is the only member of the family inside the
six-vertex catalogue.

\paragraph{The rooted factor was already in the project.}
The function $g(x,y)$ of the edge self-amalgamation proof --- the paw density
with its distinguished edge deleted --- is, for the paw,
\[
 g(x,y)=\int_\Omega W(x,z)W(y,z)d(z)\,d\mu(z),
\]
and that is exactly \texttt{Link.pageOp W 1}, built earlier for the page-rooted
book-leaf family.  So both identities the proof needs are existing lemmas:
\eqref{eq:tF-edge-rooted} is \texttt{Link.integral\_edge\_pageOp} at $s=1$
(\texttt{SelfAmalgam.integral\_edge\_page\_eq}), and
\eqref{eq:edge-amalgam-density} is one six-coordinate peeling
(\texttt{SelfAmalgam.homDensity\_amalgam}).  Nothing about amalgamation as such
is formalized; the two densities are computed directly.

\paragraph{The departure: no measure $\rho$, and no hypothesis $p>0$ on the
analytic step.}
At $N=2$ Jensen against the probability measure
$\,d\rho=W\,d\mu^2/p$ \emph{is} Cauchy--Schwarz, so the formalization never
constructs $\rho$.  It applies the project's own
\texttt{PureChordal.integral\_mul\_sq\_le\_integral\_mul\_integral\_mul\_sq} to
the \emph{unnormalised} weight $W$ on $\mu\otimes\mu$, giving
\[
 \left(\int_{\Omega^2}Wg\right)^2
 \leq\left(\int_{\Omega^2}W\right)\left(\int_{\Omega^2}Wg^2\right)
 =p\;t(A_{115},W)
 \qquad(\texttt{SelfAmalgam.sq\_edge\_page\_le}),
\]
with no hypothesis on $p$ at all.  The division by $p$ happens once, at the very
end, where $p\geq1/2>0$ is already available.  Consequently the note's standing
assumption $p=t(K_2,W)>0$ in the edge self-amalgamation lemma is not needed in
the form used here.

\paragraph{The chromatic polynomial does not come from the amalgamation
identity.}
The note derives
$\chi_{F_{115}}=\chi_{F_{15}}^2/\bigl(x(x-1)\bigr)$ from the general identity
for $\chi_{B_N^{(2)}}$.  The formalization instead builds
$\chi_{F_{115}}(x)=x(x-1)^3(x-2)^2$ from an \texttt{attachVertex} tower, the way
every other catalogue row builds its chromatic polynomial: $K_3$ on $\{0,1,2\}$,
then the leaf $3$ on $\{2\}$, then the second apex $4$ on the glued edge
$\{0,1\}$, then the leaf $5$ on $\{4\}$.  The root list $[0,1,1,1,2,2]$ then
gives $\Phi=p^3(2p-1)^2$ through \texttt{chromaticTarget\_affineProd}, with no
per-row polynomial computation --- and it is the same root list as Atlas 119.

\paragraph{The paw input.}
$t(F_{15},W)\geq p^2(2p-1)$ is
\texttt{RootedTriangleTree.rootedTree\_bound} at $r=1$, applied to
\texttt{Chromatic.pawGraph} through \texttt{paw\_factorization}.  That lemma is
stated for $\Omega:\mathrm{Type}$ rather than $\mathrm{Type}^*$, which is why
this file fixes $\Omega:\mathrm{Type}$ throughout; \texttt{SatisfiesLowerBound}
quantifies over $\mathrm{Type}$ as well, so nothing is lost.

\end{document}
